% ---
% id: reductions/rem:master-scope
% kind: rem
% title: "Branch labels as information the adversary may hold"
% summary: >-
%   The branching rule files the branch label as a block index of the system
%   handed to the adversary, so it models branches Eve may know. A setting the
%   measuring party draws privately belongs to the common-channel rule of [[lem:linear]] (ii),
%   which opens no block and produces no constraint.
% prominence: standard
% status: accepted
% uses: [lem:master, lem:linear]
% assumes: []
% refs: []
% provenance: newly-derived
% lean:
% ---

\begin{remark}\label{rem:master-scope}
The branch label ends up as a block index of $\Qsys$, and $\Qsys$ is the system
the squashed model hands to the adversary. Lemma~\ref{lem:master} therefore
models a branch Eve may as well know: a detector-response mode she selected, a
test on the light she sent, a register that arrived with it. A setting the measuring
party draws \emph{privately} can be pushed through the rule, taking
$\Ins_j=q_j\idmap$, but the resulting model files that setting in Eve's state,
and only part~(iv) of Lemma~\ref{lem:master} in its operator form ties it
back. Settings meant to stay on the measuring party's side belong to
Lemma~\ref{lem:linear}\,(ii) instead: one common channel, no block opened, no
constraint produced. The probe setting therefore enters this contract
through Lemma~\ref{lem:linear}\,(ii), and a mode register that arrived with
the light through Lemma~\ref{lem:master}; swapping the two would be a
modelling error, giving Eve a register she does not have or withholding one
she does.
\end{remark}
